\chapter{Executive Summary}

\section{System Identity and Scope}
\textbf{BengaluruTrafficAI} is an integrated, real-time intelligent transportation system (ITS) prototype developed to address urban gridlock and intersection bottlenecks. Centered on the high-density \textbf{Silk Board Junction} in Bengaluru, India, the system applies computer-vision inference, stateful vehicle tracking, proportional signal timing optimization, heuristic machine-learning classification, and commuter-facing geospatial telemetry.

The platform is designed around a dual-execution paradigm:
\begin{enumerate}[leftmargin=*]
    \item \textbf{Local Computer-Vision Mode (\texttt{backend/app.py})}: A full edge-perception server utilizing an Ultralytics YOLOv8n convolutional neural network, a pure-Python Intersection-over-Union (IoU) tracker, frame-geometry lane assignment, and a multi-type traffic violation detection suite operating on sample roadway video footage.
    \item \textbf{Cloud Demonstration Mode (\texttt{backend/app\_demo.py})}: A lightweight, micro-footprint simulation engine engineered to operate stably within constrained hosting environments (such as Render's 512\,MB RAM free tier), broadcasting unified JSON REST contracts that replicate dynamic traffic surges, signal updates, and commuter advice.
\end{enumerate}

\section{Core Capabilities}
\begin{itemize}[leftmargin=*]
    \item \textbf{Adaptive Proportional Signal Timing}: Green-light allocations are calculated dynamically between bounded limits (10 to 60 seconds) in proportion to instantaneous queue densities across the North, South, East, and West corridors.
    \item \textbf{Emergency Vehicle Preemption}: Detects prioritized emergency responders (ambulances and fire engines) and instantly overrides normal rotation to allocate an immediate 60-second green clearance corridor.
    \item \textbf{Perception-Based Infraction Auditing}: Analyzes vehicle bounding boxes to detect stop-line red-light jumps, directional wrong-side driving, rider helmet compliance, and frame-rate bounded speed estimates.
    \item \textbf{Exploratory Congestion Forecasting}: A scikit-learn Random Forest model trained on synthetic diurnal density patterns to classify intersection congestion into discrete tiers (\textit{Low}, \textit{Medium}, \textit{High}).
    \item \textbf{Commuter Advisory and Geospatial Telemetry}: A dedicated portal calculating Haversine straight-line aerial proximity from browser geolocation or landmark presets, accompanied by transparent queue delay disclosures.
    \item \textbf{Cyberpunk Operations Dashboard}: A responsive React 18 single-page application (SPA) featuring real-time junction radar visualizations, live violation event terminals, and Recharts historical trends.
\end{itemize}

\section{Honest Telemetry Disclosures and Verified Limitations}
In accordance with professional engineering integrity, this report explicitly delineates verified implementation boundaries from theoretical assumptions:
\begin{itemize}[leftmargin=*]
    \item \textbf{Uncalibrated Speed Estimations}: The speed estimation algorithm in \texttt{ai/speed\_estimator.py} calculates pixel displacement over monotonic time using an uncalibrated heuristic ratio ($0.05\,\text{m/px}$). In the absence of physical ground-plane homography rectification, outputs represent relative visual motion estimates rather than legally certified radar velocities.
    \item \textbf{Synthetic Congestion Rules}: The Random Forest model is trained on rule-generated synthetic matrices rather than empirical multi-month inductive loop sensor data. It acts as an architectural proof-of-concept for ML pipeline integration.
    \item \textbf{Optical Field-of-View Limitations}: Physical intersection cameras cannot observe queue spillover beyond their optical line of sight; clearance delay is therefore reported transparently as \texttt{ESTIMATE\_UNAVAILABLE}.
\end{itemize}
